\documentclass[10pt]{beamer}
\usetheme{Madrid}
\usepackage{booktabs}
\usepackage{graphicx}
\usepackage[T1]{fontenc}
\setbeamertemplate{navigation symbols}{}

% ======================= EDIT YOUR DETAILS =======================
\newcommand{\studentname}{Aflah Muhammed P}
% Change the university roll/registration number:
\newcommand{\rollno}{TVE23CSXXX}
% Change your class roll number:
\newcommand{\rollclass}{Roll No: XX}
% Change your guide name:
\newcommand{\guidename}{Prof. Guide Name}
% Change department / college / short-institute if needed:
\newcommand{\deptname}{Department of Computer Science and Engineering}
\newcommand{\collegename}{College of Engineering Trivandrum}
\newcommand{\shortinst}{CET}
% Change the seminar date:
\newcommand{\seminardate}{\today}
% =================================================================

\title[B.Tech Seminar]{MARS: Peer Review for AI - Making Multi-Agent Reasoning Twice as Efficient}
\subtitle{Based on: MARS: toward more efficient multi-agent collaboration for LLM reasoning (arXiv:2509.20502, 2025)}
\author[\studentname\ (\shortinst)]{\studentname \\ \small \rollno \\ \small \rollclass \\ \small Guide: \guidename}
\institute[]{\deptname \\ \collegename}
\date{\seminardate}

\begin{document}

\begin{frame}[plain]
  \titlepage
\end{frame}

\begin{frame}{Table of Contents}
  \tableofcontents
\end{frame}

\section{Introduction}
\begin{frame}{Introduction}
\begin{itemize}
  \item LLMs reason well alone but still make multi-step errors
  \item Multi-agent teamwork improves reasoning accuracy
  \item Debate-style collaboration is accurate but costly
  \item MARS reorganizes agents like academic peer review
  \item Goal: keep accuracy, cut tokens and time in half
\end{itemize}
\end{frame}

\section{Literature Survey}
\begin{frame}{Literature Survey / Background}
  \small
  \begin{tabular}{@{}p{2.7cm}p{2.3cm}p{5.6cm}@{}}
    \toprule
    \textbf{Title} & \textbf{Author} & \textbf{Summary} \\
    \midrule
    Chain-of-Thought Prompting & Wei et al., 2022 & Prompts an LLM to show intermediate reasoning steps, improving accuracy on math and logic tasks over direct answering. \\
    \addlinespace
    Multi-Agent Debate (MAD) & Du et al., 2023 & Multiple LLM agents debate over several rounds, exchanging and refining answers to improve factuality and reasoning. \\
    \addlinespace
    Reflexion / Self-Reflection & Shinn et al., 2023 & An agent critiques and revises its own output using feedback, improving results without extra models. \\
    \addlinespace
    \bottomrule
  \end{tabular}
\end{frame}

\section{Motivation}
\begin{frame}{Motivation}
\begin{itemize}
  \item Debate needs many agents and many rounds
  \item Every agent reads every other agent's text
  \item Token and time costs grow quickly
  \item Hard to deploy expensive reasoning at scale
\end{itemize}
\end{frame}

\section{Problem Statement}
\begin{frame}{Problem Statement}
  \begin{block}{}
  Multi-Agent Debate improves LLM reasoning, but its heavy agent-to-agent communication makes token usage and inference time too costly. The challenge is to keep debate-level accuracy while sharply cutting that communication overhead.
  \end{block}
\end{frame}

\section{Objectives}
\begin{frame}{Objectives}
\begin{itemize}
  \item Design a cheaper multi-agent collaboration structure
  \item Preserve the accuracy gains of debate
  \item Remove costly reviewer-to-reviewer communication
  \item Cut token usage and inference time substantially
  \item Validate across multiple benchmarks and LLMs
\end{itemize}
\end{frame}

\section{Methodology}
\begin{frame}[allowframebreaks]{Methodology / Approach}
\begin{itemize}
  \item Model collaboration on academic peer review
  \item Author agent drafts an initial solution
  \item Reviewer agents critique it independently in parallel
  \item Meta-reviewer synthesizes reviews into one verdict
  \item Author revises using meta-reviewer guidance
  \item Repeat for limited rounds, then output final answer
\end{itemize}
\end{frame}

\section{Key Concepts}
\begin{frame}[allowframebreaks]{Key Concepts}
  \begin{description}
    \item[Role-based collaboration] Agents hold fixed roles (author, reviewer, meta-reviewer) instead of symmetric debaters, structuring the workflow.
    \item[Independent review] Reviewers critique the solution without seeing each other, removing cross-talk while keeping diverse feedback.
    \item[Meta-review synthesis] A single agent merges all reviews into a final decision and concrete revision guidance.
    \item[Iterative revision] The author rewrites the solution using feedback across a small, bounded number of review rounds.
  \end{description}
\end{frame}

\section{System Architecture}
\begin{frame}{System Architecture / How It Works}
  \begin{block}{Overview}
  The pipeline is a directed flow, not a mesh. A question enters the author agent, which produces a draft solution; that draft fans out to several reviewer agents that work in parallel and independently, each returning a decision and comments. Those reviews feed into a single meta-reviewer that outputs a final decision plus revision instructions, which loop back to the author for another round if needed. A diagram should show one author node, a parallel row of reviewer nodes with no links between them, and one meta-reviewer node that closes the loop back to the author.
  \end{block}
  % TIP: insert a diagram here, e.g.:
  % \begin{center}\includegraphics[width=0.7\textwidth]{your-figure.png}\end{center}
\end{frame}

\section{Results}
\begin{frame}[allowframebreaks]{Results / Findings}
\begin{itemize}
  \item Matches MAD accuracy across MMLU, GPQA, and GSM8K
  \item Cuts token usage by roughly 50 percent
  \item Cuts inference time by roughly 50 percent
  \item Example: GPQA GPT-4o-mini about 48.3 percent vs MAD 47.5 percent
  \item Example tokens: about 7,900 for MARS vs about 17,000 for MAD
\end{itemize}
\end{frame}

\section{Discussion}
\begin{frame}{Discussion}
\begin{itemize}
  \item Communication structure drives cost as much as agent count
  \item Independent review avoids agents amplifying shared errors
  \item Clear roles make the system easier to audit and control
  \item Efficiency makes multi-agent reasoning more deployable
\end{itemize}
\end{frame}

\section{Challenges}
\begin{frame}{Limitations \& Challenges}
\begin{itemize}
  \item Still more costly than a single-agent call
  \item Reviewer feedback may overlap or be redundant
  \item Tested on selected benchmarks and model sizes
  \item Bounded rounds may cap gains on very hard problems
\end{itemize}
\end{frame}

\section{Future Work}
\begin{frame}{Future Work}
\begin{itemize}
  \item Adapt the number of reviewers or rounds dynamically
  \item Test on broader tasks and larger models
  \item Combine with tool use and retrieval
  \item Explore smarter meta-review synthesis strategies
\end{itemize}
\end{frame}

\section{Conclusion}
\begin{frame}{Conclusion}
\begin{itemize}
  \item Peer review is an efficient collaboration blueprint
  \item MARS matches debate accuracy at about half the cost
  \item Removing reviewer cross-talk is the key saving
  \item A practical step toward affordable multi-agent reasoning
\end{itemize}
\end{frame}

\section{References}
\begin{frame}[allowframebreaks]{References}
  \footnotesize
  \begin{enumerate}
    \item MARS: toward more efficient multi-agent collaboration for LLM reasoning. Xiao Wang, Jia Wang, Yijie Wang, Pengtao Dang, Sha Cao, Chi Zhang. arXiv:2509.20502, 2025.
    \item Improving Factuality and Reasoning in Language Models through Multiagent Debate. Yilun Du, Shuang Li, Antonio Torralba, Joshua B. Tenenbaum, Igor Mordatch. arXiv:2305.14325, 2023.
    \item Chain-of-Thought Prompting Elicits Reasoning in Large Language Models. Jason Wei, Xuezhi Wang, Dale Schuurmans, et al. NeurIPS, 2022.
    \item Reflexion: Language Agents with Verbal Reinforcement Learning. Noah Shinn, Federico Cassano, Ashwin Gopinath, et al. NeurIPS, 2023.
    \item Tree of Thoughts: Deliberate Problem Solving with Large Language Models. Shunyu Yao, Dian Yu, Jeffrey Zhao, et al. NeurIPS, 2023.
    \item Self-Consistency Improves Chain of Thought Reasoning in Language Models. Xuezhi Wang, Jason Wei, Dale Schuurmans, et al. ICLR, 2023.
  \end{enumerate}
\end{frame}

\begin{frame}[plain]
  \begin{center}
    {\Huge Thank You}\\[1em]
    {\large Questions?}
  \end{center}
\end{frame}

\end{document}